\section{Concrete Semantics}
\label{sec:concrete-semantics}

% TODO: ignore alignment details
% TODO: say that we assume: pointer type = i64 / index type = i32
% TODO: say the LLVM allows some freedom in the repr of memory objects / address spaces / etc.
% TODO: UB is interpreted as an error

%The most common depiction of program memory (a.k.a.\@ the ``flat model'') is a single,
%large block of memory bytes, with addresses being simple byte indexes.

In this section, we define the operational small-step semantics of \llvm IR based on the syntax in~\Cref{figure:llvm-syntax}.
The \llvm reference~\cite{llvmref} leaves some aspects of the memory model unspecified,
such as the representation of allocated objects and address spaces.
This flexibility allows different semantic interpretations,
and in this work,
we adopt the \emph{memory-object model} which is widely used in \SE engines~\cite{klee,mossberg2019manticore}.

\subsection{Memory Modeling}
\label{sec:memory-modeling}

We first define the major components of the memory model: memory objects, the heap, and the memory allocator.

\begin{definition}
\label{def:memory}
A memory object is defined as a tuple:
\[ \mobj{\addr}{\size}{\arrayvalue} \in \MemoryObjects \quad \quad \MemoryObjects \triangleq \pointertype \times \indextype \times (\Z \rightarrow \Vals^{\bvsort{8}}) \]
where $\addr$ is its base address, $\size$ is its size, and $\arrayvalue$ is a mapping that tracks the stored values,
which can be either bit-vectors of type $\bvsort{8}$ (\ie a \emph{byte}) or one of the special values $\tundef$ or $\tpoison$.
A heap $\memory \in \Heaps$ is defined as a set of memory objects.
\end{definition}

Now, we define the operations provided by the memory model:
allocation, deallocation, read, and write.

\paragraph{Allocation}
The \emph{allocation} operation is defined using the following function:
\[ \allocate : \Heaps \times \indextype \to \Heaps \times \pointertype \]
The function $\allocate$ receives a heap $\memory$ and an allocation size $\size$,
and attempts to provide a memory object $\mobj{\addr}{\size}{\arrayvalue}$ that satisfies the following:
\begin{enumerate}
    \item $[\addr, \addr + \size)$ is fully contained within the address space (\ie $\addr < \addr + \size$)
    \item $[\addr, \addr + \size)$ does not overlap with address intervals of existing memory objects
    \item $s \geq 0$ and does not exceed the largest signed integer representable by the index type
\end{enumerate}
If such a memory object can be provided, then the result is the tuple $\langle \memory \cup \mobj{\addr}{\size}{\arrayvalue}, \addr \rangle$.
Otherwise, the result is $\none$.
The mapping $\arrayvalue$ of the newly allocated memory object is set to a function that maps all the indices to $\bvval{0}{8}$.

\paragraph{Deallocation}
The \emph{deallocation} operation is defined using the following function:
\[ \deallocate : \Heaps \times \MemoryObjects \to \Heaps \]
The function $\deallocate$ receives a heap $\memory$ and a memory object $\mo$,
and returns $\memory \setminus \{\mo\}$.

\paragraph{Read}
The \emph{read} operation is defined using the following function:
\[ \moreadfunc{\tau} : \MemoryObjects \times \pointertype \to \Vals^{\tau} \]
The $\moreadfunc{\tau}$ function receives a memory object $\mo$ and a pointer $p$.\footnote{
We assume that $\tau$ is one of the built-in integer types: $\bvsort{8}$, $\bvsort{16}$, $\bvsort{32}$, $\bvsort{64}$.
}
First, it reads the consecutive $|\tau|$ cells starting from $p$ into a sequence of values of type $\Vals^{\bvsort{8}}$.
If one of these values is $\tpoison$, then the result is $\tpoison$,
and otherwise, if one of these values is $\tundef$, then the result is $\tundef$.
Otherwise, we deserialize this sequence of values (bytes) into an integer of type $\tau$.

\paragraph{Write}
The \emph{write} operation is defined using the following function:
\[ \mowritefunc{\tau} : \MemoryObjects \times \pointertype \times \Vals^{\tau} \to \MemoryObjects \]
The $\mowritefunc{\tau}$ function receives a memory object $\mo$ (\ie $\mobj{\addr}{\size}{\arrayvalue}$), a pointer $p$, and a value $v$.
First, it serializes the value $v$ into a sequence of bytes.
%Then, it uses these bytes to update the corresponding cells of the array $\arrayvalue$ starting from $p$.
Then, it uses these bytes to create a new memory object $\mobj{\addr}{\size}{\arrayvalue'}$,
where $\arrayvalue'$ is derived from $\arrayvalue$ by updating the corresponding cells of $\arrayvalue$ starting from $p$.

\paragraph{Memory Allocation Scheme}
In this paper,
we assume a linear memory allocator that allocates memory objects consecutively while maintaining a fixed gap between them,
commonly referred to as a \emph{red-zone}.
Modeling a generic memory allocator that allows arbitrary locations for memory objects would be far more challenging.
First, modeling such memory allocators requires developing a more complex theory of memory (\Cref{sec:symbolic-semantics,sec:overview,sec:program-safety,sec:proof-generation}).
Second, existing \SE engines do not support such memory allocators,
so generating safety proofs using these tools would require significant modifications to their implementations (\Cref{sec:proof-generation}).
Therefore, we leave the modeling of such memory allocators to future work.

\subsection{Operational Small-Step Semantics}

\begin{figure}
\[
\begin{array}{c}
\ruledefnonamesmall
{
    n \in \Z
}
{
    \evalexprty{n}{\localstore}{\tau} = \valoftype{n}{\tau}
}
\quad\quad
\ruledefnonamesmall
{
    \mapaccess{\localstore}{x} \in \Vals^{\tau}
}
{
    \evalexprty{x}{\localstore}{\tau} = \mapaccess{\localstore}{x}
}
\\[5pt]
\ruledefnonamesmall
{
    \evalexprty{e_1}{\localstore}{\tau} \in \Vals^{\tau}, \
    \evalexprty{e_2}{\localstore}{\tau} \in \Vals^{\tau}, \
    \sem{\toper^{\tau}} : \Vals^{\tau} \times \Vals^{\tau} \to \Vals^{\tau'}
}
{
    \evalexprty{\toper^\tau~e_1~e_2}{\localstore}{\tau'} = \sem{\toper^\tau}\big(\evalexprty{e_1}{\localstore}{\tau}, \evalexprty{e_2}{\localstore}{\tau}\big)
}
\quad\quad
\ruledefnonamesmall
{
    \evalexprty{e}{\localstore}{\tau} \in \Vals^{\tau}, \
    \sem{\tcast^{\tau\to\tau'}} : \Vals^{\tau} \to \Vals^{\tau'}
}
{
    \evalexprty{\tcast^{\tau\to\tau'}~e}{\localstore}{\tau'} = \sem{\tcast^{\tau\to\tau'}}\big(\evalexprty{e}{\localstore}{\tau}\big)
}
\end{array}
\]
\caption{Evaluation semantics of \llvm expressions.}
\label{figure:expr-semantics}
\end{figure}

\begin{figure}
\[
\begin{array}{cl}
\ruledefsmall{Branch-True}
{
  \instrat{\location} = \codequote{\tbr ~ e ~ b_1 ~ b_2} \quad
  \evalexprty{e}{\localstore}{\bvsort{1}} = \bvval{1}{1}
}
{
  \progstate{\location}{\localstore}{\stack}{\memory}
  \step
  \progstate{\instratlabel{\location}{b_1}}{\localstore}{\stack}{\memory}
}
\\[1em]
\ruledefsmall{Branch-False}
{
  \instrat{\location} = \codequote{\tbr ~ e ~ b_1 ~ b_2} \quad
  \evalexprty{e}{\localstore}{\bvsort{1}} = \bvval{0}{1}
}
{
  \progstate{\location}{\localstore}{\stack}{\memory}
  \step
  \progstate{\instratlabel{\location}{b_2}}{\localstore}{\stack}{\memory}
}
\\[1em]
\ruledefsmall{Make-Symbolic}{
  \instrat{\location} = \codequote{x ~\texttt{=}~ \tcall^{\bvsort{32}} ~ \intrinsicfunc{make\_symbolic}()} \quad
  n \in \Nat
}
{
  \progstate{\location}{\localstore}{\stack}{\memory}
  \step
  \progstate{\nextinstr{\location}}{\localstore[x \mapsto \bvval{n}{32}]}{\stack}{\memory}
}
\\[1em]
\ruledefsmall{Malloc}{
  \instrat{\location} = \codequote{x~\texttt{=}~\tcall^{\pointertypeof{\bvsort{8}}}~\intrinsicfunc{malloc}(\bvsort{64} ~ e)} \quad
  \evalexprty{e}{\localstore}{\bvsort{64}} = \bvval{n}{64} \quad
  \allocate(\memory, \valoftype{n}{\indextype}) = (\memory', p)
}
{
  \progstate{\location}{\localstore}{\stack}{\memory}
  \step
  \progstate{\nextinstr{\location}}{\localstore[x \mapsto p]}{\stack}{\memory'}
}
\\[1em]
\ruledefsmall{Free}{
  \instrat{\location} = \codequote{\tcall~\intrinsicfunc{free}(\pointertypeof{\bvsort{8}} ~ e)} \quad
  \evalexprty{e}{\localstore}{\pointertypeof{\bvsort{8}}} = p \quad
  \mo \in \memory \quad
  \isbaseaddr{p}{\mo} \quad
  \deallocate(\memory, \mo) = \memory'
}
{
  \progstate{\location}{\localstore}{\stack}{\memory}
  \step
  \progstate{\nextinstr{\location}}{\localstore}{\stack}{\memory'}
}
\\[1em]
\ruledefsmall{Load}{
  \instrat{\location} = \codequote{x~\texttt{=}~\tload^{\tau}~e} \quad
  \evalexprty{e}{\localstore}{\pointertypeof{\tau}} = p \quad
  \mo \in \memory \quad
  \isinbounds{p}{\tau}{\mo}
}
{
  \progstate{\location}{\localstore}{\stack}{\memory}
  \step
  \progstate{\nextinstr{\location}}
    {\localstore[x \mapsto \moread{\tau}{\mo}{p}]}
    {\stack}{\memory}
}
\\[1em]
\ruledefsmall{Store}{
  \instrat{\location} = \codequote{\tstore^{\tau}~e_1~e_2} \quad
  \evalexprty{e_1}{\localstore}{\tau} = v \quad
  \evalexprty{e_2}{\localstore}{\pointertypeof{\tau}} = p \quad
  \mo \in \memory \quad
  \isinbounds{p}{\tau}{\mo}
}
{
  \progstate{\location}{\localstore}{\stack}{\memory}
  \step
  \progstate{\nextinstr{\location}}{\localstore}{\stack}{\heapreplace{\memory}{\mo}{\mowrite{\tau}{\mo}{p}{v}}}
}
\\[1em]
\ruledefsmall{GEP}{
  \instrat{\location} = \codequote{x~\texttt{=}~\tgep~\pointertypeof{\tau}~e~\overline{(\tau_i, e_i)}} \quad
  \evalexprty{e}{\localstore}{\pointertypeof{\tau}} = p \quad
  \evalexprty{e_i}{\localstore}{\tau_i} = v_i
}
{
  \progstate{\location}{\localstore}{\stack}{\memory}
  \step
  \progstate{\nextinstr{\location}}
    {\localstore[x \mapsto \evalgep(\tau, p, \overline{v_i})]}
    {\stack}{\memory}
}
\\[2em]
\scalebox{0.8}{$
\begin{array}{l@{\quad}l}
\module & \mbox{module} \\
\instrat{\location} & \mbox{instruction at location $\location$ in $\module$} \\
\nextinstr{\location} & \mbox{next location} \\
\instratlabel{\location}{b} & \mbox{location at label $b$} \\
\isbaseaddr{p}{\mo} & \mbox{$p$ is equal to the base address of $\mo$} \\
\isinbounds{p}{\tau}{\mo} & \mbox{interval $[p, p + \sizeof{\tau})$ is contained within the address interval of $\mo$} \\
\moread{\tau}{\mo}{p} & \mbox{read a value of type $\tau$ from memory object $\mo$ starting from address $p$} \\
\mowrite{\tau}{\mo}{p}{v} & \mbox{write value $v$ into memory object $\mo$ at address $p$} \\
\heapreplace{\memory}{\mo}{\mo'} & \mbox{update $\memory$ by replacing memory object $\mo$ with $mo'$}
\end{array}
$}
\end{array}
\]
\caption{Operational small-step semantics of \llvm IR.}
\label{figure:llvm-semantics}
\end{figure}

In the following,
we define the concrete state and the derivation rules of the operational small-step semantics.
In this paper,
we consider a concrete state with a local store and a heap.
However, our framework (in \coq) and implementation (on top of \klee) support also stack frames, stack allocations (\ie $\talloca$), and global variables.

% TODO: define a frame
\begin{definition}
\label{def:state}
%A location $\location\in\Locs$ is specified by a function identifier, block identifier, instruction index,
%and an optional previous-block identifier.
A \emph{concrete state} is a tuple $\progstate{\location}{\localstore}{\stack}{\memory}$,
where $\location \in \Locs$ is a program location,
$\localstore \in \Stores$ is a \emph{local store} that maps variables to values in $\Vals$,
and $\memory \in \Heaps$ is a heap.
\end{definition}

\Cref{figure:expr-semantics} defines the evaluation semantics of \llvm expressions.
The typed semantics of an expression, denoted by $\evalexprty{e}{\localstore}{\tau}$,
defines the value of $e$ given a local store $\localstore \in \Stores$ in a context where a value of type $\tau$ is expected.
Operators are parameterized by bit-vector types,
and their semantics are functions whose domains and ranges depend on those types (\eg $\sem{\textsf{add}^\tau} : \Vals^\tau\times\Vals^\tau\to\Vals^\tau$).

\Cref{figure:llvm-semantics} defines the operational small-step semantics for a subset of the instructions introduced in~\Cref{figure:llvm-syntax}.
In the following, we discuss the derivation rules in more detail.\footnote{
The full definition can be found in our \coq framework.
}

\textbf{\textit{Branch.}}
The instruction $\tbr$ follows standard operational semantics.
We first evaluate the branch condition $e$
and then jump to the corresponding location based on the resulting value.

\textbf{\textit{Make-Symbolic.}}
The \intrinsicfunc{make\_symbolic} function is an intrinsic function provided by the symbolic execution environment to introduce fresh symbolic values.
In the concrete semantics,
we model this behavior by non-deterministically assigning the variable $x$ to an arbitrary concrere value $\bvval{n}{32}$.

\textbf{\textit{Malloc.}}
When allocating a memory object,
we call the $\allocate$ function (\Cref{sec:memory-modeling}).
If $\allocate$ succeeds, then we update the local store and the heap accordingly.
Otherwise, the result is undefined behavior.

%When allocating a memory object,
%the size of the allocation must be non-negative and must not exceed the largest signed integer representable by the index type.
%As mentioned before, the address interval of the allocated memory object may not cross the unsigned boundary of the address space.
%If one of these properties is violated, the result is undefined behavior.
%The mapping of the newly allocated memory object is set to a function that maps all the indices to $\bvval{0}{8}$.

\textbf{\textit{Free.}}
% TODO: assume that the type of p is iptr?
When deallocating a pointer $p$,
we first check if there exists a memory object $\mobj{\addr}{\size}{\arrayvalue}$ in the heap $\memory$ such that $p = \addr$.
%If that is the case, then we update the heap using the $\deallocate$ function.
If that is the case, we remove that memory object from $\memory$ using the $\deallocate$ function (\Cref{sec:memory-modeling}).
Otherwise, the result is undefined behavior.

\textbf{\textit{Load and Store.}}
% TODO: what happens if p is undef/poison?
When accessing $n$ bytes from a pointer $p$,
we need to make sure that there exists a memory object $\mobj{\addr}{\size}{\arrayvalue}$ in the heap $\memory$ such that
(1) $p < p + n$ (\ie there is no overflow) and (2) $p \geq \addr \ \land \ p + n < \addr + \size$.
If that is not the case, then the result is undefined behavior.
Otherwise, we perform a read or write operation.
%For a read operation,
%we first read the corresponding sequence of stored values.
%If one of these values is $\tpoison$, we return $\tpoison$,
%and otherwise, if one of these values is $\tundef$, we return $\tundef$.
%Otherwise, we deserialize this sequence of values (bytes) into an integer of the appropriate type.
For a read operation,
we call the $\moreadfunc{\tau}$ function (\Cref{sec:memory-modeling}) to read the value from $\mo$ and update the local store accordingly.
%For a write operation,
%we serialize the value $v$ into a sequence of bytes and update the array $\arrayvalue$ accordingly.
For a write operation,
we call the $\mowritefunc{\tau}$ function (\Cref{sec:memory-modeling}) to update the contents of $\mo$ accordingly.

%\textbf{\textit{Store.}}
%% TODO: what happens if p/v is undef/poison?
%When writing a value $v$ to a pointer value $p$,
%we need to make sure that there exists a memory object $\mobj{\addr}{\size}{\arrayvalue}$ in the memory $\memory$ such that
%(1) $p < p + n$ and (2) $\addr \leq p < \addr + \size$.
%If that is not the case, then the result is undefined behavior.
%Otherwise, we use the \textit{write} procedure to store the value $v$.
%This procedure serializes the value $v$ into a sequence of bytes and updates the array $\arrayvalue$ accordingly.

\textbf{\textit{GetElementPtr.}}
In general, the $\tgep$ instruction computes the address of a subelement of an aggregate data structure (array, struct, \etc).
The function $\evalgep$ defines the semantics of the $\tgep$ instruction based on~\Cref{alg:gep} (function \textsc{gep}).\footnote{
The semantics presented here consider only arrays, however, our framework and implementation support also structs.
}
The algorithm uses the functions $\signextfunc{\indexwidth}$ and $\zeroextfunc{\indexwidth}$ defined in~\Cref{sec:llvm-ir}.
The first argument of $\evalgep$ is the type of the base pointer,
the second argument is the base pointer to start from,
and the remaining arguments are indices that indicate which of the elements of the aggregate object are indexed.
All the indices are first converted to values of the index type (sign-extended or truncated), and then translated into offsets (\cref{line:gep:array-index,line:gep:first-index}).
The first index is used to move accross the elements corresponding to the base pointer type,
so it is multiplied by the type size of the element type of the base pointer type (\cref{line:gep:first-index}).
The remaining indices are translated into offsets based on the indexed type.
If the indexed type is an array type,
then the index is multiplied by the type size of the indexed type (\cref{line:gep:array-index}).
%If the indexed type is a struct type,
%then the index is translated into a field selector.
The computed offsets are sequentially summed up,
and then the resulting offset is converted to the pointer type and added to the base pointer (\cref{line:gep:add-offset}).
%
The $\tgep$ instruction may have attributes (\opcode{nuw}, \opcode{nusw}, and \opcode{inbounds}) that impose additional rules.
For example, when the \opcode{nuw} attribute is enabled, the following conditions must hold:
(1) if the type of an index is larger than the index type, the truncation to the index type must preserve the unsigned value,
(2) the multiplication of an index by the type size does not cause an unsigned overflow of the index type, and
(3) the successive addition of each offset does not cause an unsigned overflow of the index type.
If any of these conditions is violated, the result is $\tpoison$.

\begin{algorithm}[t]
\small
\caption{Simplified semantics of $\tgep$.}
\label{alg:gep}
\begin{algorithmic}[1]

\Function{compute-offset}{$\tau, o, \overline{j}$}
    \If {$|\overline{j}| = 0$}
        \State \Return $o$
    \EndIf
    \State $j_1 :: \overline{j'} \gets \overline{j}$
    \If {$\tau = \arraytype{\tau'}{n}$}
        \State $o' \gets \signext{\indexwidth}{j_1} \cdot \valoftype{|\tau'|}{\indextype}$ \label{line:gep:array-index}
        \State \Return \Call{compute-offset}{$\tau', o + o', \overline{j'}$}
    \EndIf
\EndFunction

\Function{gep}{$\tau, p, \overline{j}$}
    \If {$\tpoison \in \overline{j}$}
        \State \Return $\tpoison$
    \EndIf
    \If {$\tundef \in \overline{j}$}
        \State \Return $\tundef$
    \EndIf
    \State $j_1 :: \overline{j'} \gets \overline{j}$
    \State $o \gets \signext{\indexwidth}{j_1} \cdot \valoftype{|\tau|}{\indextype}$ \label{line:gep:first-index}
    \State $r \gets \Call{compute-offset}{\tau, o, \overline{j'}}$
    \State \Return $p + \zeroext{\pointerwidth}{r}$ \label{line:gep:add-offset}
\EndFunction

\end{algorithmic}
\end{algorithm}

In the following,
we define when a concrete state is considered to be an error state.
We then define when a program, \ie an \llvm module, is considered safe.
Intuitively, an error arises when there is an operation that triggers undefined behavior or an operation that involves $\tpoison$ values.
The formal definition is given below.

% TODO: some text before?
% TODO: say that some notations are defined in Figure X
\begin{definition}
\label{def:error-state}
Let $c \triangleq \progstate{\location}{\localstore}{\stack}{\memory}$ be a concrete state.
Then $c$ is an error state, denoted by $\errorstate(c)$, if one of the following holds:
\begin{enumerate}
    \item $\instrat{\location} = \textup{$\codequote{\tunreachable}$}$
    \item $\instrat{\location} = \textup{$\codequote{x \ \texttt{=} \ op \ \bvsort{w} \ e_1 \ e_2}$}$, and one of the following holds (where $v_k \triangleq \evalexprty{e_k}{\localstore}{\bvsort{w}}$):
    \begin{enumerate}[left=-3pt]
        \item \textup{$op \in \{\opcode{udiv}, \opcode{sdiv}, \opcode{urem}, \opcode{srem}\}$ and $v_2 = \bvval{0}{w}$}
        \item \textup{$op = \opcode{sdiv}$, $v_1 = \bvval{(-2^{w - 1})}{w}$, and $v_2 = \bvval{(-1)}{w}$}
        \item \textup{$op \in \{\opcode{shl}, \opcode{lshr}, \opcode{ashr}\}$ and $v_2 = \bvval{n}{w}$, and $n \geq w$}
    \end{enumerate}
    \item $\instrat{\location} = \textup{$\codequote{x \ \texttt{=} \ \tcall^{\pointertypeof{\bvsort{8}}} \ \intrinsicfunc{malloc}(\bvsort{64} \ e)}$}$,
          and \textup{$\evalexprty{e}{\localstore}{\bvsort{64}} > \bvval{(2^{\indexwidth - 1} - 1)}{64}$}.
    \item $\instrat{\location} = \textup{$\codequote{\tcall~\intrinsicfunc{free}(\pointertypeof{\bvsort{8}} \ e)}$}$, \textup{$\evalexprty{e}{\localstore}{\pointertypeof{\bvsort{8}}} = p$},
          and there is no $\mo \in \memory$ s.t. $\isbaseaddr{p}{\mo}$.
    \item $\instrat{\location} = \textup{$\codequote{x \ \texttt{=} \ \tload^{\tau} \ e}$}$, \textup{$\evalexprty{e}{\localstore}{\pointertypeof{\tau}} = p$},
          and there is no $\mo \in \memory$ s.t. $\isinbounds{p}{|\tau|}{\mo}$.
    \item $\instrat{\location} = \textup{$\codequote{\tstore^{\tau} \ e_1 \ e_2}$}$, \textup{$\evalexprty{e_2}{\localstore}{\pointertypeof{\tau}} = p$},
          and there is no $\mo \in \memory$ s.t. $\isinbounds{p}{|\tau|}{\mo}$.
    % TODO: $|\overline{(\tau_i, e_i)}| > 0$
    \item $\instrat{\location} = \textup{$\codequote{x \ \texttt{=} \ \tgep~\pointertypeof{\tau} \ e \ \overline{(\tau_i, e_i)}}$}$ and $\evalgep(\tau, p, \overline{v_i}) = \tpoison$
          (where \textup{$p \triangleq \evalexprty{e}{\localstore}{\pointertypeof{\tau}}$} and $v_i \triangleq \evalexprty{e_i}{\localstore}{\tau_i}$).
\end{enumerate}
\end{definition}

\begin{definition}
\label{def:safety}
Let $\location$ be the entry program location of the module $\module$.
The \emph{initial state} of $\module$, denoted by $\initstate(\module)$,
is defined as $\progstate{\location}{\{\}}{[]}{[]}$.
Now, let $c_0 \triangleq \initstate(\module)$.
Then the module $\module$ is \emph{safe}, denoted by $\issafe(\module)$,
if for every $c$ s.t. $c_{0} \multistep c$, $c$ is not an error state (the relation $\multistep$ is the reflexive-transitive closure of $\step$).
\end{definition}
